% ============================================================================
% PART 4: THE SCIENCECLUSTER
% ============================================================================

\sectionslide{Part 4. The UZH ScienceCluster}{Three doors to the same machines}

\begin{frame}{Why a Cluster}

\small
\begin{itemize}
    \setlength{\itemsep}{5pt}
    \item \textbf{GPUs your laptop does not have}: L4 (24 GB), A100 (80 GB), H100, H200. A 27-billion-parameter model on an L4; a 70B on an A100 or H100.
    \item \textbf{Long jobs}: a whole corpus annotated overnight, while your laptop is closed.
    \item \textbf{Data that must stay at UZH}: confidential texts, licensed data. The model comes to the data.
\end{itemize}

\vspace{0.3cm}

\textbf{Access}
\begin{itemize}
    \setlength{\itemsep}{3pt}
    \item Your group needs a Science IT project; your PI adds you. Use is billed to the project.
    \item Off campus: the UZH VPN first.
    \item \textbf{Two rules.} Login nodes are for editing and submitting, never for computing. Large files (model weights) go on \texttt{/scratch}, which is purged after 30 days unused. Nothing is backed up.
\end{itemize}

\vspace{0.2cm}
{\footnotesize Every command of this part: \texttt{setup/cluster\_guide.md} in the course repository.}

\end{frame}

\begin{frame}[fragile]{Door 1. The Terminal}

\small
\textbf{Log in, look around}
{\scriptsize
\begin{verbatim}
ssh -l <shortname> cluster.s3it.uzh.ch     # UZH network or VPN
quota                                       # your space
module available                            # the software
\end{verbatim}}

\textbf{A GPU, for an hour, interactively}
{\scriptsize
\begin{verbatim}
module load l4
srun --pty --time=1:0:0 --mem-per-cpu=7800 --cpus-per-task=2 --gpus=1 bash -l
nvidia-smi                                  # you are on a compute node now
exit                                        # GPU released
\end{verbatim}}

\textbf{A job, while you do something else}
{\scriptsize
\begin{verbatim}
sbatch annotate.sh       # a script with #SBATCH lines: time, memory, --gpus=1
squeue --me              # PD = waiting, R = running
\end{verbatim}}

\end{frame}

\begin{frame}[fragile]{Door 2. An Agent on the Cluster}

\small
\textbf{The same agent, in a folder on the cluster}
{\scriptsize
\begin{verbatim}
curl -fsSL https://claude.ai/install.sh | bash   # once, on the login node
claude                          # press c, open the URL on your laptop, paste the code
\end{verbatim}}

\vspace{0.1cm}

\textbf{What it is good at there}: writing job scripts, submitting them, reading the logs, relaunching what failed. The well-defined task that can run unattended (Session 2).

\vspace{0.2cm}

\textbf{Three precautions}
\begin{itemize}
    \setlength{\itemsep}{3pt}
    \item \textbf{One line in \texttt{CLAUDE.md}}: \textit{``Never run analysis code on the login node: write a job script and submit it with sbatch.''}
    \item \textbf{Your home is shared by all nodes}: turn self-updates off (\texttt{DISABLE\_UPDATES}), or a session on another node can crash.
    \item \textbf{VS Code or Cursor}: S3IT asks you not to connect them to the login nodes. Its documentation has an advanced recipe through a job.
\end{itemize}

\end{frame}

\begin{frame}{Door 3. The Browser: ScienceApps}

\small
\centering
{\large \href{https://apps.s3it.uzh.ch}{\texttt{apps.s3it.uzh.ch}}}\\[0.1cm]
{\footnotesize UZH login, private window. Every cluster account has access.}

\vspace{0.3cm}
\raggedright
\begin{itemize}
    \setlength{\itemsep}{4pt}
    \item \textbf{Jupyter, RStudio}, on a compute node, with the GPU you choose in a form: type, hours (up to a week), cores, memory.
    \item \textbf{A full desktop}, on GPU nodes too.
    \item \textbf{VS Code in the browser} (beta), without its AI features.
    \item \textbf{Files}: upload, download, and a shell, in the browser.
\end{itemize}

\vspace{0.3cm}

\textbf{The simplest door if you never used a terminal}: a Jupyter session with one L4, today's notebook, and Part 5 against a model served on that node.

\end{frame}

\begin{frame}[fragile]{A Local Model on a GPU Node}

\small
Inside a GPU session. Apptainer runs the official Ollama image; S3IT supports Apptainer but has no page on LLMs, so this is the course's recipe.

{\scriptsize
\begin{verbatim}
module load apptainer
export OLLAMA_MODELS=/scratch/$USER/ollama-models     # weights on scratch
apptainer pull /scratch/$USER/ollama.sif docker://ollama/ollama
apptainer exec --nv /scratch/$USER/ollama.sif ollama serve > ollama.log 2>&1 &
apptainer exec --nv /scratch/$USER/ollama.sif ollama pull gemma3:27b
\end{verbatim}}

\vspace{0.1cm}

\begin{itemize}
    \setlength{\itemsep}{3pt}
    \item The notebook's Part 5 then works unchanged, on that node: Ollama answers on \texttt{localhost:11434}.
    \item \textbf{A whole corpus}: the same lines in a batch job. \textbf{Many requests in parallel}: vLLM, an OpenAI-compatible server, with \texttt{make\_openai\_annotator}.
    \item Do not open the port to others (UZH IT's advice).
\end{itemize}

\vspace{0.15cm}
{\footnotesize Outside UZH: the CSCS Inference API serves open models (Apertus, Qwen, Gemma) to groups with a CSCS project, through S3IT.}

\end{frame}
